% Comandos para crear la portada y la primera página del PFG
%% Título del PFG
\newcommand{\PFGtitulo}{Título del PFG}
%% Título corto, para mostrar en el encabezado de las páginas pares. Por defecto se pone el título, pero, si es muy largo, puede acortarse aquí
\newcommand{\PFGtitulocorto}{\PFGtitulo}
%% Autor/a y directores/as del PFG
\newcommand{\PFGautor}{Nombre del autor/a del PFG}
\newcommand{\PFGdirector}{Director/a del PFG}
%% Descomentar y rellenar la siguiente línea si hay codirector/a en el PFG
%\newcommand{\PFGcodirector}{Codirector/a del PFG}
%% Grado. A elegir entre:
%% Ingeniería Informática
%% Ingeniería en Tecnologías de la Información
\newcommand{\PFGgrado}{Grado cursado}
%% Modalidad. A elegir entre:
%% genérica
%% específica
\newcommand{\PFGmodalidad}{modalidad}
%% Curso académico en formato 20xx/20xx+1. Por ejemplo 2022/2023
\newcommand{\PFGcursoacademico}{Curso académico}
%% Convocatoria. A elegir entre:
%% junio
%% septiembre
%% diciembre
\newcommand{\PFGconvocatoria}{Convocatoria}
%% Fecha de lectura (si no se sabe el día concreto, poner mes y año}
\newcommand{\PFGfechadefensa}{Fecha de lectura}


\documentclass[12pt,twoside]{book}
\usepackage{graphicx}
\usepackage[utf8]{inputenc}
\usepackage[papersize={210mm,297mm},tmargin=25mm,bmargin=25mm,lmargin=20mm,rmargin=20mm]{geometry}

% No indentar los párrafos y separarlos 1cm
\setlength\parindent{0cm}
\setlength\parskip{1em}

% Para tener los elementos de texto en español
\usepackage[spanish]{babel}
% Para poder resaltar un texto en amarillo con \hl{}
\usepackage{soul}
% Definición de colores
\usepackage{xcolor}
\definecolor{links}{rgb}{0.35, 0.4, 0.35}
\definecolor{urls}{rgb}{0,0,1}
\definecolor{cites}{rgb}{0.8,0,0}
\definecolor{gray20}{gray}{0.20}

% Hiperreferencias
\usepackage[hyphens]{url}
\usepackage{hyperref}
\hypersetup{
	colorlinks = true,
	linkcolor = links,
	urlcolor = urls,
	citecolor = cites
}


% Formato de los listados
\usepackage{listings}
\lstloadlanguages{C,C++,Haskell,HTML,Java,PHP,Python}

\definecolor{codegreen}{rgb}{0,0.6,0}
\definecolor{codegray}{rgb}{0.5,0.5,0.5}
\definecolor{codepurple}{rgb}{0.58,0,0.82}
\definecolor{backcolour}{rgb}{0.95,0.95,0.92}

\lstdefinestyle{mystyle}{
    commentstyle=\color{codegreen},
    keywordstyle=\color{blue},
    numberstyle=\tiny\color{codegray},
    stringstyle=\color{codepurple},
    backgroundcolor=\color{backcolour}, % choose the background color; you must add \usepackage{color} or \usepackage{xcolor}; should come as last argument
	basicstyle=\footnotesize\ttfamily\color{black},        % the size of the fonts that are used for the code
	belowcaptionskip=1\baselineskip,
	breakatwhitespace=false,         % sets if automatic breaks should only happen at whitespace
	breaklines=true,                 % sets automatic line breaking
	captionpos=b,                    % sets the caption-position to bottom
	deletekeywords={...},            % if you want to delete keywords from the given language
	escapeinside={\%*}{*)},          % if you want to add LaTeX within your code
	extendedchars=true,              % lets you use non-ASCII characters; for 8-bits encodings only, does not work with UTF-8
	firstnumber=1,                % start line enumeration with line 1000
	frame=single,	                   
	%identifierstyle=\color{blue!50!black},
	keepspaces=true,                 % keeps spaces in text, useful for keeping indentation of code (possibly needs columns=flexible)
	%language=C,                 % the language of the code
	%morekeywords={*,...},            % if you want to add more keywords to the set
	numbers=left,                    % where to put the line-numbers; possible values are (none, left, right)
	numbersep=10pt,                   % how far the line-numbers are from the code
	numberstyle=\tiny\color{gray}, % the style that is used for the line-numbers
	rulecolor=\color{black},         % if not set, the frame-color may be changed on line-breaks within not-black text (e.g. comments (green here))
	morecomment=[n]{(*}{*)},
	morecomment=[n][\color{magenta}]{\%I}{\ },
	morecomment=[n][\color{magenta}]{\%Q}{\ },
	showspaces=false,                % show spaces everywhere adding particular underscores; it overrides 'showstringspaces'
	showstringspaces=false,          % underline spaces within strings only
	showtabs=false,                  % show tabs within strings adding particular underscores
	stepnumber=1,                    % the step between two line-numbers. If it's 1, each line will be numbered
	tabsize=2,	                   % sets default tabsize to 2 spaces
	title=\lstname,
	xleftmargin=1.5cm,
}
\lstset{style=mystyle}




% Añadir ficheros externos
\usepackage{subfiles}
% Soporte para anexos
\usepackage{appendix}
% Paquetes para figuras
\usepackage{caption}
\usepackage{subcaption}
% Para las tablas
\usepackage{tabularx}
\usepackage{multicol}
\usepackage{multirow}
%Encabezados bonitos y acordes al reglamento
\usepackage{fancyhdr}
\fancyhf{}
\lhead[\thepage]{\leftmark}
\rhead[\PFGtitulocorto]{\thepage}
\setlength{\headheight}{16pt}
\pagestyle{empty}

% Paquete para poner la bibliografia en la table of contents
\usepackage[nottoc,notlot,notlof]{tocbibind}
% Cambiamos el espacio en blanco enorme que por defecto hay en las páginas inicio de capítulo
\usepackage{titlesec}
\titleformat{\chapter}[display]   
{\normalfont\huge\bfseries}{\chaptertitlename\ \thechapter}{20pt}{\Huge}   
\titlespacing*{\chapter}{0pt}{-50pt}{40pt}

% Para crear subsubsubsection
\setcounter{secnumdepth}{4}
\titleformat{\paragraph}
{\normalfont\normalsize\bfseries}{\theparagraph}{1em}{}
\titlespacing*{\paragraph}
{0pt}{3.25ex plus 1ex minus .2ex}{1.5ex plus .2ex}

% Ruta en la que pondremos las imágenes para el documento
\graphicspath{ {imagenes/} }
% Redifinimos el nombre que asigna babel spanish a las tablas (cuadro) a tabla
\renewcommand{\spanishtablename}{Tabla}
\renewcommand{\spanishlisttablename}{Índice de tablas}
\renewcommand{\spanishcontentsname}{Índice}
% Redefinimos el nombre de apéndice a anexo
\renewcommand{\appendixname}{Anexos}
\renewcommand{\appendixtocname}{Anexos}
\renewcommand{\appendixpagename}{Anexos}

% Definimos el comando de página en blanco
\newcommand\paginablanco{%
    \null
    \thispagestyle{empty}%
    \newpage}

% Definimos el comando de página en blanco sin avanzar numeracion
\newcommand\paginablancosin{%
    \paginablanco{}
    \addtocounter{page}{-1}}

% Mayor anchura en las tablas
\renewcommand{\arraystretch}{1.5}
% Mayor ancuhura en las fracciones
\newcommand\ddfrac[2]{\frac{\displaystyle #1}{\displaystyle #2}}
% Comando escribir PFG
\newcommand{\pfg}{Proyecto Fin de Grado }


% Comenzamos el documento
\begin{document}
\pagenumbering{gobble} % para no incluir números en la portada y primera página

% Hacemos la portada
\begin{titlepage}

\begin{figure}[!htb]
    \centering
    \includegraphics{rosadelosvientos.png}
\end{figure}

\begin{center}
    {\scshape\Large Universidad Nacional de Educación a Distancia \par}
    {\scshape\Large Escuela Técnica Superior de Ingeniería Informática \par}
    \vfill
    \large{Proyecto de fin de Grado en \PFGgrado}\\
    \vfill
    {\bfseries\Huge \textcolor{black}{\PFGtitulo} \par}
    \vfill
\end{center}

\vspace*{3cm}
{\huge \textcolor{black}{\PFGautor} \par}
{\huge \textcolor{black}{Dirigido por: \PFGdirector} \par}
\ifdefined\PFGcodirector
  \huge \textcolor{black}{Codirigido por: \PFGcodirector} \par
\fi
{\huge \textcolor{black}{Curso \PFGcursoacademico, convocatoria \PFGconvocatoria} \par}

\chapter*{}

\thispagestyle{empty}

\begin{figure}[!htb]
    \centering
    \includegraphics{rosadelosvientos.png}
\end{figure}

\begin{center}

    {\bfseries\Huge \textcolor{black}{\PFGtitulo} \par}
    \vfill
    {\bfseries\Large \textcolor{black}{Proyecto de fin de Grado en \PFGgrado} \par}
    {\bfseries\Large \textcolor{black}{de modalidad \PFGmodalidad} \par}
\end{center}

\vspace*{5cm}
{\Large Realizado por: \PFGautor \par}
{\Large Dirigido por: \PFGdirector \par}
\ifdefined\PFGcodirector
  {\Large Codirigido por: \PFGcodirector \par}
\fi%

%%{\Large Codirigido por: \PFGcodirector \par}
\vfill
{\Large Fecha de lectura y defensa: \PFGfechadefensa \par}

\end{titlepage}

\pagenumbering{Roman} % para comenzar la numeracion de paginas en numeros romanos

% Agradecimientos
\chapter*{Agradecimientos}
\thispagestyle{empty}
\begin{flushright}
\textit{Incluir todos los agradecimientos}
\end{flushright}

\chapter*{Resumen} % El * indica que no queremos que añada la palabra "Capitulo"

Resumen del PFG.

\chapter*{Abstract} % El * indica que no queremos que añada la palabra "Capitulo"

Resumen del PFG en inglés.

\chapter*{Palabras clave} % El * indica que no queremos que añada la palabra "Capitulo"

Palabras clave del PFG

\chapter*{Keywords} % El * indica que no queremos que añada la palabra "Capitulo"

Palabras clave del PFG en inglés.

\tableofcontents %% Índice general

%% SI NO HAY ÍNDICE DE TABLAS, COMENTAR LAS DOS SIGUIENTES LÍNEAS
\listoftables    %% Índice de tablas
\addcontentsline{toc}{section}{Índice de tablas} % Queremos que aparezca en el índice

%% SI NO HAY ÍNDICE DE FIGURAS, COMENTAR LAS DOS SIGUIENTES LÍNEAS
\listoffigures   %% Índice de figuras
\addcontentsline{toc}{section}{Índice de figuras} % Queremos que aparezca en el índice
\newpage

\paginablanco{}

%% Comenzamos con la numeración arábiga y el estilo de cabeceras de páginas definido, porque aquí empieza el contenido
\pagenumbering{arabic}
\pagestyle{fancy}
%% Ejemplo de importación de capítulos
\subfile{capitulos/C1-introduccion}
\subfile{capitulos/C2-estado_arte}
\subfile{capitulos/C3-propuesta}
\subfile{capitulos/C4-disenno}
\subfile{capitulos/C5-desarrollo}
\subfile{capitulos/C6-pruebas}
\subfile{capitulos/C7-conclusiones}

%% Bibliografía
\newpage
\subfile{capitulos/referencias}

\clearpage

%% Anexos
\appendix
\subfile{capitulos/A1-anexo}
\subfile{capitulos/A2-anexo}

%...

\end{document}